% Source: https://artofproblemsolving.com/community/c3h2808290_2012_geometry_bee__preliminary_round
% Stable identity: aops-topic-2808290-q08
Each bounding circle has radius 8, because a 60-degree chord equals its radius. Their centres are $4\sqrt3$ from the chord midpoint, on opposite sides. The lens has total thickness $2(8-4\sqrt3)$ perpendicular to the chord. Any inscribed disk has diameter at most this thickness. The disk centred at the chord midpoint with radius $8-4\sqrt3$ lies inside both bounding disks, by the triangle inequality for distances to their centres. It attains the bound, so $a=8,b=-4,c=3$ and $\boxed{a+b+c=7}$.
\begin{asy}
usepackage("amsmath");
size(7cm);
pair A=(-4,0),B=(4,0),U=(0,4*sqrt(3)),V=(0,-4*sqrt(3));draw(arc(U,8,240,300));draw(arc(V,8,60,120));draw(circle((0,0),8-4*sqrt(3)),blue);draw(A--B,dashed);label("A",A,W);label("B",B,E);label("8",(0,0),S);
\end{asy}
